\chapter{Introdução}\label{chap1}
\acresetall


\section{Enquadramento e Motivação}
% Descreve o contexto geral do estágio, por que razão este projeto é importante
% e qual a motivação pessoal/profissional para a realização do mesmo.

Escreve aqui sobre o contexto geral do teu estágio.

\section{Apresentação da Entidade de Acolhimento}
% Apresenta a empresa "20 Recolher". O que faz, a sua área de atuação,
% onde se localiza e a sua importância no mercado. Podes adicionar secções
% sobre o organograma ou equipa se for relevante.

A 20 Recolher é uma empresa...

\section{Objetivos do Estágio}
% Define claramente os objetivos gerais e específicos que pretendes atingir
% ao longo do teu estágio na empresa.

Os principais objetivos definidos para este estágio incluem:
\begin{itemize}
    \item Objetivo 1;
    \item Objetivo 2;
\end{itemize}

\section{Estrutura do Relatório}
% Faz um resumo do conteúdo de cada capítulo do relatório para guiar o leitor.

O presente documento está estruturado da seguinte forma:
\begin{itemize}
    \item \textbf{Capítulo~\ref{chap1} -- Introdução:} Apresenta o enquadramento, a empresa de acolhimento, os objetivos e a organização do relatório;
    \item \textbf{Capítulo~\ref{chap2} -- Enquadramento Empresarial:} Apresenta a entidade de acolhimento (20 Recolher), a equipa de trabalho e o enquadramento do estagiário;
    \item \textbf{Capítulo~\ref{chap3} -- Atividades Planeadas:} Descreve os objetivos definidos e o plano de trabalhos inicial;
    \item \textbf{Capítulo~\ref{chap4} -- Atividades Desenvolvidas e Resultados:} Detalha as tarefas realizadas, dividindo-as entre o desenvolvimento de software (aplicações e base de dados) e o trabalho manual de hardware (testes de componentes e triagem de resíduos eletrónicos), apresentando os respetivos resultados obtidos;
    \item \textbf{Capítulo~\ref{chap5} -- Sugestões, Comentários e Propostas de Melhoria:} Reúne reflexões críticas, sugestões e propostas de evolução técnica para o projeto;
    \item \textbf{Capítulo~\ref{chap6} -- Conclusões:} Faz um balanço final dos objetivos atingidos e das competências adquiridas.
\end{itemize}
